% =============================================================================
% preambulo.tex — Formato del anexo metodológico
%
% Deliberadamente sobrio: es un documento de referencia técnica, para consultar
% y para auditar, no una pieza de divulgación. Los colores son los del sistema
% visual del proyecto (02_Code/04_figuras/DISENO.md) para que el anexo y las
% figuras del informe se lean como una sola familia.
% =============================================================================

\usepackage[utf8]{inputenc}
\usepackage[T1]{fontenc}
\usepackage[spanish,es-noquoting,es-tabla]{babel}
\usepackage{lmodern}
\usepackage{microtype}

\usepackage[a4paper,margin=2.5cm,headheight=14pt]{geometry}

% Una ruta larga sin espacios no se puede partir y TeX prefiere desbordar la
% caja antes que estirar la línea. Con este margen de emergencia estira.
\emergencystretch=3em
\usepackage{booktabs}
\usepackage{tabularx}
\usepackage{longtable}
\usepackage{array}
\usepackage{ragged2e}
\usepackage{enumitem}
\usepackage{fancyhdr}
\usepackage{titlesec}
\usepackage{xcolor}
\usepackage{graphicx}
\usepackage[most]{tcolorbox}
\usepackage{caption}
\usepackage[hidelinks]{hyperref}

% --- Colores del sistema visual (DISENO.md §2) -------------------------------
\definecolor{azulmarino}{HTML}{4458A6}
\definecolor{baya}{HTML}{C24A49}
\definecolor{verdeselva}{HTML}{159449}
\definecolor{tintauno}{HTML}{292E38}
\definecolor{tintados}{HTML}{595E66}
\definecolor{tintatres}{HTML}{888C94}
\definecolor{grilla}{HTML}{DCDEE1}
\definecolor{zafreclaro}{HTML}{A1B4DB}

\hypersetup{colorlinks=true, linkcolor=azulmarino, urlcolor=azulmarino,
            citecolor=azulmarino}

% --- Títulos ------------------------------------------------------------------
\titleformat{\section}{\normalfont\Large\bfseries\color{tintauno}}{\thesection}{1em}{}
\titleformat{\subsection}{\normalfont\large\bfseries\color{tintauno}}{\thesubsection}{1em}{}
\titleformat{\subsubsection}{\normalfont\normalsize\bfseries\color{tintados}}{\thesubsubsection}{1em}{}
\titlespacing*{\section}{0pt}{2.2ex plus 1ex minus .2ex}{1.2ex plus .2ex}

% --- Encabezado y pie ---------------------------------------------------------
\pagestyle{fancy}
\fancyhf{}
\renewcommand{\headrulewidth}{0.4pt}
\fancyhead[L]{\footnotesize\color{tintatres} Proyecto Provincias · Anexo metodológico}
\fancyhead[R]{\footnotesize\color{tintatres}\nouppercase{\leftmark}}
\fancyfoot[C]{\footnotesize\color{tintatres}\thepage}

% --- Tablas -------------------------------------------------------------------
\renewcommand{\arraystretch}{1.15}
\captionsetup{font=small, labelfont={bf,color=tintados},
              textfont={color=tintados}, justification=raggedright,
              singlelinecheck=false}
\newcolumntype{Y}{>{\RaggedRight\arraybackslash}X}

% --- Cajas --------------------------------------------------------------------
% Tres roles, y solo tres: una regla que el pipeline aplica, un hallazgo de la
% auditoría y una advertencia sobre el dato. Más tipos de caja y el lector deja
% de distinguirlas.

\newtcolorbox{regla}[1][]{
  enhanced, breakable, colback=white, colframe=azulmarino,
  boxrule=0pt, leftrule=2.5pt, arc=0pt, outer arc=0pt,
  left=8pt, right=8pt, top=6pt, bottom=6pt,
  colbacktitle=white, fonttitle=\bfseries\small, coltitle=azulmarino,
  title={Regla del pipeline}, #1}

\newtcolorbox{hallazgo}[1][]{
  enhanced, breakable, colback=white, colframe=baya,
  boxrule=0pt, leftrule=2.5pt, arc=0pt, outer arc=0pt,
  left=8pt, right=8pt, top=6pt, bottom=6pt,
  colbacktitle=white, fonttitle=\bfseries\small, coltitle=baya,
  title={Hallazgo de la auditoría}, #1}

\newtcolorbox{advertencia}[1][]{
  enhanced, breakable, colback=grilla!30, colframe=grilla,
  boxrule=0.4pt, arc=1pt,
  left=8pt, right=8pt, top=6pt, bottom=6pt,
  colbacktitle=grilla!30, fonttitle=\bfseries\small, coltitle=tintados,
  title={Advertencia sobre el dato}, #1}

% --- Diagramas ----------------------------------------------------------------
% Los estilos del DAG se declaran aquí una sola vez. Cada tipo de nodo
% corresponde a una etapa del flujo y cada tipo de arista a una forma de
% dependencia; los .tex generados solo referencian estos nombres, nunca colores
% ni tamaños sueltos. La leyenda del documento (§2.2) describe exactamente esta
% tabla de estilos: si aquí cambia un color, allí cambia la leyenda.

\usepackage{tikz}
\usetikzlibrary{arrows.meta, positioning, calc, fit, backgrounds, shapes.geometric}

\tikzset{
  % --- Nodos, uno por etapa del flujo ---
  nodobase/.style={
    draw, rounded corners=1.5pt, line width=0.5pt,
    font=\scriptsize, inner sep=2.6pt, align=left,
    minimum height=4.6mm, anchor=west},
  nodofuente/.style={nodobase, draw=tintatres, fill=grilla!35, text=tintauno},
  nodoscript/.style={nodobase, draw=azulmarino, fill=azulmarino!8, text=tintauno},
  nododerivado/.style={nodobase, draw=zafreclaro, fill=zafreclaro!25, text=tintauno},
  nodorecurso/.style={nodobase, draw=tintatres, fill=white, text=tintauno},
  nodoseccion/.style={nodobase, draw=azulmarino, fill=azulmarino!15, text=tintauno},
  nodosalida/.style={nodobase, draw=verdeselva, fill=verdeselva!10, text=tintauno},
  % --- Aristas ---
  flujo/.style={-{Stealth[length=1.4mm, width=1.1mm]},
                draw=tintatres, line width=0.28pt, opacity=0.75},
  flujoexcepcion/.style={-{Stealth[length=1.4mm, width=1.1mm]},
                draw=baya, line width=0.34pt, densely dashed, opacity=0.9},
  % --- Etapas del diagrama de arquitectura ---
  % `text width` fijo: sin él, una etiqueta larga ensancha la caja hasta
  % solaparse con la vecina y las flechas quedan sin sitio donde dibujarse.
  cajabase/.style={rounded corners=2pt, line width=0.6pt, align=center,
                font=\footnotesize, text width=2.9cm, inner sep=4pt,
                minimum height=1.15cm, text=tintauno},
  etapa/.style={cajabase, draw=azulmarino, fill=azulmarino!8},
  almacen/.style={cajabase, draw=tintatres, fill=grilla!40},
  salida/.style={cajabase, draw=verdeselva, fill=verdeselva!10},
  paso/.style={-{Stealth[length=2mm, width=1.6mm]}, draw=azulmarino,
               line width=0.8pt},
  control/.style={-{Stealth[length=2mm, width=1.6mm]}, draw=baya,
               line width=0.7pt, densely dashed},
}

% --- Atajos -------------------------------------------------------------------
\newcommand{\ruta}[1]{\texttt{\small #1}}
\newcommand{\fn}[1]{\texttt{\small #1}}
% Muestra de estilo para la leyenda de §2.2: dibuja el propio nodo o la propia
% arista, para que la leyenda no pueda desincronizarse de lo que declara.
\newcommand{\muestranodo}[1]{%
  \raisebox{-0.6mm}{\begin{tikzpicture}[baseline]
    \node[#1, minimum width=9mm, minimum height=3.4mm] at (0,0) {};
  \end{tikzpicture}}}
\newcommand{\muestraarista}[1]{%
  \raisebox{0.4mm}{\begin{tikzpicture}[baseline]
    \draw[#1] (0,0) -- (0.9,0);
  \end{tikzpicture}}}
